\documentclass[a4paper]{article}     
\usepackage{amsfonts}
\usepackage{amsmath}
\usepackage{amssymb}
\usepackage{graphicx}
\DeclareMathOperator{\cis}{cis}

\begin{document}
	
	\noindent \large Auteur: Victor Pena \\
	\noindent \large Studierichting: Informatica\\
	\noindent \large Oefeningenreeks 1D nr. 21.8\\
	
	\medskip
	
	\normalsize 
	
	$P(z)=z^3 -z^2 + 3z - 10$ \\
	
	
	$P(2)=8 - 4 + 6 - 10 = 0$\\
	
	$\Rightarrow$ \begin{tabular}{l|lll|l}
		 & 1 & -1 & 3  &-10 \\
		2&   & 2  & 2  &10 \\  \cline{1-5}
		 & 1 & 1  & 5  &0
	\end{tabular} $\Rightarrow P(z) = (z-2)(z^2 + z + 5)$\\
	
	
	De wortels van $z^2 + z + 5$ zijn gegeven door: \\
	
	$z = \dfrac{-1 \pm i\sqrt{20-1}}{2} = \dfrac{-1 \pm i\sqrt{19}}{2} = -\dfrac{1}{2} \pm i\dfrac{\sqrt{19}}{2}$\\
	
	
	Dus, $P(z) = (z-2)(z+\dfrac{1}{2}-i\dfrac{\sqrt{19}}{2})(z+\dfrac{1}{2}+i\dfrac{\sqrt{19}}{2})$, met wortels $2,-\dfrac{1}{2}+i\dfrac{\sqrt{19}}{2}$ en $-\dfrac{1}{2}-i\dfrac{\sqrt{19}}{2}$ \\
	
	
	
	
	
	
	
	
	
	\begin{thebibliography}{999}
		%\bibitem{a} Referentie
	\end{thebibliography}
\end{document} 
